% Options for packages loaded elsewhere
\PassOptionsToPackage{unicode}{hyperref}
\PassOptionsToPackage{hyphens}{url}
%
\documentclass[
  12pt,
]{article}
\usepackage{amsmath,amssymb}
\usepackage{iftex}
\ifPDFTeX
  \usepackage[T1]{fontenc}
  \usepackage[utf8]{inputenc}
  \usepackage{textcomp} % provide euro and other symbols
\else % if luatex or xetex
  \usepackage{unicode-math} % this also loads fontspec
  \defaultfontfeatures{Scale=MatchLowercase}
  \defaultfontfeatures[\rmfamily]{Ligatures=TeX,Scale=1}
\fi
\usepackage{lmodern}
\ifPDFTeX\else
  % xetex/luatex font selection
\fi
% Use upquote if available, for straight quotes in verbatim environments
\IfFileExists{upquote.sty}{\usepackage{upquote}}{}
\IfFileExists{microtype.sty}{% use microtype if available
  \usepackage[]{microtype}
  \UseMicrotypeSet[protrusion]{basicmath} % disable protrusion for tt fonts
}{}
\makeatletter
\@ifundefined{KOMAClassName}{% if non-KOMA class
  \IfFileExists{parskip.sty}{%
    \usepackage{parskip}
  }{% else
    \setlength{\parindent}{0pt}
    \setlength{\parskip}{6pt plus 2pt minus 1pt}}
}{% if KOMA class
  \KOMAoptions{parskip=half}}
\makeatother
\usepackage{xcolor}
\usepackage[margin=2.5cm]{geometry}
\usepackage{color}
\usepackage{fancyvrb}
\newcommand{\VerbBar}{|}
\newcommand{\VERB}{\Verb[commandchars=\\\{\}]}
\DefineVerbatimEnvironment{Highlighting}{Verbatim}{commandchars=\\\{\}}
% Add ',fontsize=\small' for more characters per line
\usepackage{framed}
\definecolor{shadecolor}{RGB}{248,248,248}
\newenvironment{Shaded}{\begin{snugshade}}{\end{snugshade}}
\newcommand{\AlertTok}[1]{\textcolor[rgb]{0.94,0.16,0.16}{#1}}
\newcommand{\AnnotationTok}[1]{\textcolor[rgb]{0.56,0.35,0.01}{\textbf{\textit{#1}}}}
\newcommand{\AttributeTok}[1]{\textcolor[rgb]{0.13,0.29,0.53}{#1}}
\newcommand{\BaseNTok}[1]{\textcolor[rgb]{0.00,0.00,0.81}{#1}}
\newcommand{\BuiltInTok}[1]{#1}
\newcommand{\CharTok}[1]{\textcolor[rgb]{0.31,0.60,0.02}{#1}}
\newcommand{\CommentTok}[1]{\textcolor[rgb]{0.56,0.35,0.01}{\textit{#1}}}
\newcommand{\CommentVarTok}[1]{\textcolor[rgb]{0.56,0.35,0.01}{\textbf{\textit{#1}}}}
\newcommand{\ConstantTok}[1]{\textcolor[rgb]{0.56,0.35,0.01}{#1}}
\newcommand{\ControlFlowTok}[1]{\textcolor[rgb]{0.13,0.29,0.53}{\textbf{#1}}}
\newcommand{\DataTypeTok}[1]{\textcolor[rgb]{0.13,0.29,0.53}{#1}}
\newcommand{\DecValTok}[1]{\textcolor[rgb]{0.00,0.00,0.81}{#1}}
\newcommand{\DocumentationTok}[1]{\textcolor[rgb]{0.56,0.35,0.01}{\textbf{\textit{#1}}}}
\newcommand{\ErrorTok}[1]{\textcolor[rgb]{0.64,0.00,0.00}{\textbf{#1}}}
\newcommand{\ExtensionTok}[1]{#1}
\newcommand{\FloatTok}[1]{\textcolor[rgb]{0.00,0.00,0.81}{#1}}
\newcommand{\FunctionTok}[1]{\textcolor[rgb]{0.13,0.29,0.53}{\textbf{#1}}}
\newcommand{\ImportTok}[1]{#1}
\newcommand{\InformationTok}[1]{\textcolor[rgb]{0.56,0.35,0.01}{\textbf{\textit{#1}}}}
\newcommand{\KeywordTok}[1]{\textcolor[rgb]{0.13,0.29,0.53}{\textbf{#1}}}
\newcommand{\NormalTok}[1]{#1}
\newcommand{\OperatorTok}[1]{\textcolor[rgb]{0.81,0.36,0.00}{\textbf{#1}}}
\newcommand{\OtherTok}[1]{\textcolor[rgb]{0.56,0.35,0.01}{#1}}
\newcommand{\PreprocessorTok}[1]{\textcolor[rgb]{0.56,0.35,0.01}{\textit{#1}}}
\newcommand{\RegionMarkerTok}[1]{#1}
\newcommand{\SpecialCharTok}[1]{\textcolor[rgb]{0.81,0.36,0.00}{\textbf{#1}}}
\newcommand{\SpecialStringTok}[1]{\textcolor[rgb]{0.31,0.60,0.02}{#1}}
\newcommand{\StringTok}[1]{\textcolor[rgb]{0.31,0.60,0.02}{#1}}
\newcommand{\VariableTok}[1]{\textcolor[rgb]{0.00,0.00,0.00}{#1}}
\newcommand{\VerbatimStringTok}[1]{\textcolor[rgb]{0.31,0.60,0.02}{#1}}
\newcommand{\WarningTok}[1]{\textcolor[rgb]{0.56,0.35,0.01}{\textbf{\textit{#1}}}}
\usepackage{longtable,booktabs,array}
\usepackage{calc} % for calculating minipage widths
% Correct order of tables after \paragraph or \subparagraph
\usepackage{etoolbox}
\makeatletter
\patchcmd\longtable{\par}{\if@noskipsec\mbox{}\fi\par}{}{}
\makeatother
% Allow footnotes in longtable head/foot
\IfFileExists{footnotehyper.sty}{\usepackage{footnotehyper}}{\usepackage{footnote}}
\makesavenoteenv{longtable}
\usepackage{graphicx}
\makeatletter
\def\maxwidth{\ifdim\Gin@nat@width>\linewidth\linewidth\else\Gin@nat@width\fi}
\def\maxheight{\ifdim\Gin@nat@height>\textheight\textheight\else\Gin@nat@height\fi}
\makeatother
% Scale images if necessary, so that they will not overflow the page
% margins by default, and it is still possible to overwrite the defaults
% using explicit options in \includegraphics[width, height, ...]{}
\setkeys{Gin}{width=\maxwidth,height=\maxheight,keepaspectratio}
% Set default figure placement to htbp
\makeatletter
\def\fps@figure{htbp}
\makeatother
\setlength{\emergencystretch}{3em} % prevent overfull lines
\providecommand{\tightlist}{%
  \setlength{\itemsep}{0pt}\setlength{\parskip}{0pt}}
\setcounter{secnumdepth}{-\maxdimen} % remove section numbering
\usepackage{setspace}
\usepackage{amsmath, amsfonts, amssymb}
\usepackage{booktabs}
\usepackage{csquotes}
\usepackage{graphicx}
\usepackage{xcolor}
\usepackage{hyperref}
\usepackage{etoolbox}
\AtBeginEnvironment{Shaded}{\setstretch{1}}
\usepackage{titlesec}
\titlespacing*{\section}{0pt}{*0.7}{*0.3}
\titlespacing*{\subsection}{0pt}{*0.5}{*0.2}
\titlespacing*{\subsubsection}{0pt}{*0.3}{*0.1}
\ifLuaTeX
  \usepackage{selnolig}  % disable illegal ligatures
\fi
\usepackage[style=numeric,]{biblatex}
\addbibresource{quellen.bib}
\nocite{R-manual_2021}
\usepackage{bookmark}
\IfFileExists{xurl.sty}{\usepackage{xurl}}{} % add URL line breaks if available
\urlstyle{same}
\hypersetup{
  pdftitle={Multilevel models for longitudinal data},
  pdfauthor={Maria Levina, Rahul Vishwkarma},
  hidelinks,
  pdfcreator={LaTeX via pandoc}}

\title{Multilevel models for longitudinal data}
\author{Maria Levina, Rahul Vishwkarma}
\date{2026-03-31}

\begin{document}
\maketitle

{
\setcounter{tocdepth}{2}
\tableofcontents
}
\newpage

\section{Motivation}\label{motivation}

Longitudinal studies are very common in various fields like psychology,
medicine, education, economics, and sociology, and numerous examples and
collections of longitudinal data can be found online (see, for instance,
\href{https://atlaslongitudinaldatasets.ac.uk/}{Atlas of longitudinal
studies}). Longitudinal studies are designed to observe a sample over
time to extract knowledge about patterns of trends and fluctuations in
the general population. The same group of subjects, which may be single
persons or groups/organizations, is tracked over time and its
characteristics of interest are observed.

In this work, we will take a closer look at an alternative method to
analyze longitudinal data using the multilevel modeling framework. This
method has several advantages over the traditional approaches and allows
analyzing more real-world longitudinal datasets. Here we will use an
example dataset to showcase the multilevel modeling approach and go
through the analysis step by step, providing R code to build models and
create figures and its output.

\section{Setup}\label{setup}

\subsection{Import}\label{import}

In the following we will use implementation of multilevel models from
the package \texttt{lme4} \autocite{lme4_2015} as well as statistical
tests from package \texttt{lmerTest} \autocite{lmertest_2017}. We mostly
use \texttt{ggplot2} \autocite{ggplot2_2016} for graphics and
\texttt{knitr} and \texttt{pander} \autocite{knitr_2014} for R Markdown
generation.

\begin{Shaded}
\begin{Highlighting}[]
  \FunctionTok{library}\NormalTok{(}\StringTok{"lme4"}\NormalTok{)}
  \FunctionTok{library}\NormalTok{(}\StringTok{"lmerTest"}\NormalTok{)}
  \FunctionTok{library}\NormalTok{(}\StringTok{"ggplot2"}\NormalTok{)}
  \FunctionTok{library}\NormalTok{(}\StringTok{"dplyr"}\NormalTok{)}
\end{Highlighting}
\end{Shaded}

\subsection{Example dataset}

As mentioned above we will use an example to illustrate multilevel
approach. Here a dataset from the collection of Centre for Multilevel
Modelling, University of Bristol is taken, which can be accessed via
\href{https://www.bristol.ac.uk/cmm/learning/support/datasets/}{this
link}. We have chosen a dataset describing weight gains of 568 Asian
children in a British community over several weeks after birth. The
children were measured on up to five occasions visiting a clinic roughly
to target examination dates of 6 weeks, then 8, 12, 27 months.
Additional variables recorded were weight at birth and sex of a baby.

The code below shows how to load the dataset and convert it to format
which suits best for further analysis.

\begin{Shaded}
\begin{Highlighting}[]
  \CommentTok{\# Read the file line by line and convert to one single string}
\NormalTok{  data\_path }\OtherTok{=} \StringTok{"asian/ASIAN.DAT"}
\NormalTok{  data\_raw }\OtherTok{\textless{}{-}} \FunctionTok{toString}\NormalTok{(}\FunctionTok{readLines}\NormalTok{(data\_path))}

  \CommentTok{\# Split the long string by commas}
\NormalTok{  data\_raw }\OtherTok{\textless{}{-}} \FunctionTok{strsplit}\NormalTok{(data\_raw, }\StringTok{", "}\NormalTok{)}

  \CommentTok{\# Create a one{-}column data frame}
\NormalTok{  dataset }\OtherTok{\textless{}{-}} \FunctionTok{data.frame}\NormalTok{(data\_raw)}
  \FunctionTok{colnames}\NormalTok{(dataset) }\OtherTok{\textless{}{-}} \StringTok{"raw"}

  \CommentTok{\# Extract columns}
\NormalTok{  extract\_columns }\OtherTok{\textless{}{-}} \ControlFlowTok{function}\NormalTok{(idxs, df)\{}
\NormalTok{      column }\OtherTok{=} \FunctionTok{sapply}\NormalTok{(df, substr, }\AttributeTok{start =}\NormalTok{ idxs[}\DecValTok{1}\NormalTok{], }\AttributeTok{stop =}\NormalTok{ idxs[}\DecValTok{2}\NormalTok{])}
      \FunctionTok{return}\NormalTok{(column)\}}

  \CommentTok{\# Define start and end positions for each column}
\NormalTok{  col\_idxs }\OtherTok{=} \FunctionTok{list}\NormalTok{(}\FunctionTok{c}\NormalTok{(}\DecValTok{1}\NormalTok{,}\DecValTok{4}\NormalTok{), }\CommentTok{\# ChildID: characters 1{-}4}
                  \FunctionTok{c}\NormalTok{(}\DecValTok{5}\NormalTok{,}\DecValTok{7}\NormalTok{), }\CommentTok{\# Age: characters 5{-}7}
                  \FunctionTok{c}\NormalTok{(}\DecValTok{8}\NormalTok{,}\DecValTok{12}\NormalTok{), }\CommentTok{\# Weight: characters 8{-}12}
                  \FunctionTok{c}\NormalTok{(}\DecValTok{13}\NormalTok{,}\DecValTok{16}\NormalTok{), }\CommentTok{\# Birthweight: characters 13{-}16}
                  \FunctionTok{c}\NormalTok{(}\DecValTok{17}\NormalTok{,}\DecValTok{17}\NormalTok{)) }\CommentTok{\# Gender: character 17}

\NormalTok{  col\_list }\OtherTok{\textless{}{-}} \FunctionTok{lapply}\NormalTok{(col\_idxs, extract\_columns, }\AttributeTok{df =}\NormalTok{ dataset) }\CommentTok{\# 5}

  \CommentTok{\# Assign extracted columns to new variables in the data frame}
\NormalTok{  dataset}\SpecialCharTok{$}\NormalTok{ChildID }\OtherTok{\textless{}{-}} \FunctionTok{as.vector}\NormalTok{(col\_list[[}\DecValTok{1}\NormalTok{]])}
\NormalTok{  dataset}\SpecialCharTok{$}\NormalTok{Age }\OtherTok{\textless{}{-}} \FunctionTok{as.vector}\NormalTok{(col\_list[[}\DecValTok{2}\NormalTok{]])}
\NormalTok{  dataset}\SpecialCharTok{$}\NormalTok{Weight }\OtherTok{\textless{}{-}} \FunctionTok{as.vector}\NormalTok{(col\_list[[}\DecValTok{3}\NormalTok{]])}
\NormalTok{  dataset}\SpecialCharTok{$}\NormalTok{Birthweight }\OtherTok{\textless{}{-}} \FunctionTok{as.vector}\NormalTok{(col\_list[[}\DecValTok{4}\NormalTok{]])}
\NormalTok{  dataset}\SpecialCharTok{$}\NormalTok{Gender }\OtherTok{\textless{}{-}} \FunctionTok{as.vector}\NormalTok{(col\_list[[}\DecValTok{5}\NormalTok{]])}

\NormalTok{  dataset }\OtherTok{\textless{}{-}}\NormalTok{ dataset[,}\SpecialCharTok{{-}}\DecValTok{1}\NormalTok{]}
\end{Highlighting}
\end{Shaded}

\section{Challenges of longitudinal
data}\label{challenges-of-longitudinal-data}

\subsection{Levels of analysis and time
continuum}\label{levels-of-analysis-and-time-continuum}

Longitudinal studies are very popular despite being quite complex and
expensive to convey \autocite{hoffman_2015}. It naturally depends on the
scope of a study: how large the study sample is, how long are study
subjects tracked for and even how probable a drop-out is. However,
people are hard to track in general and measures need to be
well-documented over time and comparable, i.e.~measure methods need to
stay consistent.

Nevertheless, analysis of repeated measures of several subjects is quite
useful. They make different levels of analysis accessible:

\begin{enumerate}
\def\labelenumi{\arabic{enumi}.}
\item
  Within-person / intra-individual analysis: relationships among
  intra-individual differences in time-variant variables. (level 1
  analysis)
\item
  Between-person / inter-individual analysis: relationships among
  inter-individual differences in time-invariant variables. (level 2
  analysis)
\end{enumerate}

It is also important to differentiate, whether change or fluctuation in
target value should be investigated. In \textcite{hoffman_2015} the two
terms are use to describe different kinds of longitudinal variability.
Here, change describes systematic increase or decrease over time, often
in longer perspective, while fluctuation usually refers to undirected
variation over time on short-term horizons.

\subsection{Model for mean and model for
variance}\label{model-for-mean-and-model-for-variance}

A model for mean describes structural part of the model, i.e.~how does
the outcome depend on regressor values for a study subject. It includes
\emph{fixed effects} of the covariates, where \emph{fixed} means that
the effect is same for all subjects. In the model for the mean a choice
should be made, what regressors should be taken in account and in which
form. This choice is usually based on specific research question.

A model for variance describes, on the other hand, the distribution of
residuals of the outcome. Separate model for variance allows us to model
patterns of variance and covariance in residuals, which introduces great
flexibility compared to overly simplifying assumptions for variance
structure in generalized linear models or ANOVA. There residuals are
often assumed to be normally distributed and independent across persons,
with constant variance across persons. That often doesn't hold for
real-world longitudinal data \autocite{hoffman_2015}. Modeling variance
separately also allows to include so-called \emph{random effects} of
covariates, when every study subject has its own effect of the same
predictor. Addition of random effects allow greater flexibility -
variance between and within a person can be distinguished.

\subsection{Models for repeated
measures}\label{models-for-repeated-measures}

Traditional approach for modeling longitudinal data - \textbf{R}epeated
\textbf{M}easures \textbf{An}alysis \textbf{o}f \textbf{Va}riance
(RM-ANOVA) accounts for dependency of residuals within one unit/subject,
but still assumes normal distribution of residuals as well as
homoscedasticity (constant variance over time points) and compound
symmetry (equal variances and equal covariances over subjects).

Moreover, the goal of RM-ANOVA model is not to represent the change and
fluctuation in the target variable in some general functional form, but
to ``reproduce'' the observed mean per occasion using fixed effects
equal to the number of occasions minus 1 \autocite[p.96]{hoffman_2015}.
This sets an additional constraint to data such as equal intervals
between observation time points as well as same time points for every
participant, which is often hard to implement in real life.

Also one could imagine that it may happen quite often, that a
participant doesn't show up once or drops of out of study completely due
to any possible reason. RM-ANOVA cannot handle missing data, because it
requires all observations for all subjects and time points to be
present. Therefore all those subjects where at least one observation is
missing is removed from the dataset. And many datasets cannot be
analyzed with this method at all. Multilevel models for longitudinal
data cover all of the named disadvantages and have less strict
assumptions about data distribution. Therefore, they represent an
attractive alternative to RM-ANOVA.

\subsection{Data preprocessing}

An important part of any analysis is data exploration and preparation.
In the example above, the dataset consists of five variables: ID of a
child, age (in days), weight (in grams), birth weight (in grams) and
gender (encoded as 1 or 2). First of all, we want to ensure the correct
type of variable for each column (integer/factor/string). We also change
encoding of gender to dummy encoding (with 1 for girls and 0 for boys
respectively) and store numeric values and category name in two separate
columns (see code below).

We also calculate number of observations for each child to see, whether
the dataset is balanced across subjects. We can already see from the
first 6 observations that this is not the case here.

\begin{Shaded}
\begin{Highlighting}[]
  \CommentTok{\# Preprocess columns}
\NormalTok{  dataset[, }\DecValTok{2}\SpecialCharTok{:}\DecValTok{4}\NormalTok{] }\OtherTok{\textless{}{-}} \FunctionTok{sapply}\NormalTok{(dataset[,}\DecValTok{2}\SpecialCharTok{:}\DecValTok{4}\NormalTok{], as.integer)}
\NormalTok{  dataset}\SpecialCharTok{$}\NormalTok{ChildID }\OtherTok{\textless{}{-}} \FunctionTok{as.factor}\NormalTok{(}\FunctionTok{gsub}\NormalTok{(}\StringTok{" "}\NormalTok{, }\StringTok{""}\NormalTok{, dataset}\SpecialCharTok{$}\NormalTok{ChildID))}
 
  \CommentTok{\# Decode gender}
\NormalTok{  dataset}\SpecialCharTok{$}\NormalTok{GenderID }\OtherTok{\textless{}{-}} \FunctionTok{as.integer}\NormalTok{(dataset}\SpecialCharTok{$}\NormalTok{Gender) }\SpecialCharTok{{-}} \DecValTok{1}
\NormalTok{  dataset}\SpecialCharTok{$}\NormalTok{Gender }\OtherTok{\textless{}{-}} \FunctionTok{factor}\NormalTok{(dataset}\SpecialCharTok{$}\NormalTok{GenderID, }\AttributeTok{labels =} \FunctionTok{c}\NormalTok{(}\StringTok{"Boy"}\NormalTok{, }\StringTok{"Girl"}\NormalTok{))}

  \CommentTok{\# Number of observations}
\NormalTok{  dataset}\SpecialCharTok{$}\NormalTok{NObs }\OtherTok{\textless{}{-}} \FunctionTok{as.vector}\NormalTok{(}\FunctionTok{table}\NormalTok{(dataset}\SpecialCharTok{$}\NormalTok{ChildID)[dataset}\SpecialCharTok{$}\NormalTok{Child])}
  \FunctionTok{kable}\NormalTok{(}\FunctionTok{head}\NormalTok{(dataset))}
\end{Highlighting}
\end{Shaded}

\begin{longtable}[]{@{}lrrrlrr@{}}
\toprule\noalign{}
ChildID & Age & Weight & Birthweight & Gender & GenderID & NObs \\
\midrule\noalign{}
\endhead
\bottomrule\noalign{}
\endlastfoot
2 & 377 & 10262 & 3360 & Girl & 1 & 2 \\
2 & 825 & 14500 & 3360 & Girl & 1 & 2 \\
3 & 54 & 4053 & 3310 & Girl & 1 & 4 \\
3 & 242 & 8136 & 3310 & Girl & 1 & 4 \\
3 & 377 & 9752 & 3310 & Girl & 1 & 4 \\
3 & 825 & 13180 & 3310 & Girl & 1 & 4 \\
\end{longtable}

Furthermore, we can plot observations against time to see, how much the
measurement times vary from child to child - we already know from the
dataset description, that intervals between measurements are not equal
and may differ from case to case.

\begin{Shaded}
\begin{Highlighting}[]
\NormalTok{plot\_all\_obs }\OtherTok{\textless{}{-}} \ControlFlowTok{function}\NormalTok{(varx, vary, df,}
                             \AttributeTok{transfx =} \StringTok{""}\NormalTok{, }\AttributeTok{transfy =} \StringTok{""}\NormalTok{)\{}
    \FunctionTok{ggplot}\NormalTok{(df, }\FunctionTok{aes}\NormalTok{(}\AttributeTok{x =}\NormalTok{ .data[[varx]], }\AttributeTok{y =}\NormalTok{ .data[[vary]])) }\SpecialCharTok{+} 
    \FunctionTok{coord\_cartesian}\NormalTok{(}\AttributeTok{xlim =} \FunctionTok{c}\NormalTok{(}\FunctionTok{min}\NormalTok{(df[varx]), }\FunctionTok{max}\NormalTok{(df[varx])),}
                    \AttributeTok{ylim =} \FunctionTok{c}\NormalTok{(}\FunctionTok{min}\NormalTok{(df[vary]), }\FunctionTok{max}\NormalTok{(df[vary]))) }\SpecialCharTok{+}
    \FunctionTok{geom\_line}\NormalTok{(}\FunctionTok{aes}\NormalTok{(}\AttributeTok{group=}\NormalTok{ChildID), }\AttributeTok{col =} \StringTok{"darkgray"}\NormalTok{) }\SpecialCharTok{+} 
    \FunctionTok{geom\_point}\NormalTok{(}\FunctionTok{aes}\NormalTok{(}\AttributeTok{fill=}\FunctionTok{as.factor}\NormalTok{(ChildID)), }\AttributeTok{pch=}\DecValTok{21}\NormalTok{, }\AttributeTok{size=}\DecValTok{2}\NormalTok{, }\AttributeTok{stroke=}\DecValTok{1}\NormalTok{) }\SpecialCharTok{+} 
    \FunctionTok{scale\_x\_continuous}\NormalTok{(}\AttributeTok{name =} \FunctionTok{paste0}\NormalTok{(}\StringTok{"Age"}\NormalTok{,transfx), }
                       \AttributeTok{breaks=}\FunctionTok{round}\NormalTok{(}\FunctionTok{seq}\NormalTok{(}\FunctionTok{min}\NormalTok{(df[varx]), }\FunctionTok{max}\NormalTok{(df[varx]), }\AttributeTok{length.out =} \DecValTok{10}\NormalTok{)))  }\SpecialCharTok{+} 
    \FunctionTok{scale\_y\_continuous}\NormalTok{(}\AttributeTok{name =} \FunctionTok{paste0}\NormalTok{(}\StringTok{"Weight"}\NormalTok{, transfy),}
                       \AttributeTok{limits=}\FunctionTok{range}\NormalTok{(df[,vary])) }\SpecialCharTok{+} 
    \FunctionTok{theme\_bw}\NormalTok{() }\SpecialCharTok{+} \FunctionTok{theme}\NormalTok{(}\AttributeTok{axis.text.x=}\FunctionTok{element\_text}\NormalTok{(}\AttributeTok{size=}\DecValTok{12}\NormalTok{, }\AttributeTok{colour=}\StringTok{"black"}\NormalTok{),}
                       \AttributeTok{axis.text.y=}\FunctionTok{element\_text}\NormalTok{(}\AttributeTok{size=}\DecValTok{12}\NormalTok{, }\AttributeTok{colour=}\StringTok{"black"}\NormalTok{), }
                       \AttributeTok{axis.title=}\FunctionTok{element\_text}\NormalTok{(}\AttributeTok{size=}\DecValTok{12}\NormalTok{,}\AttributeTok{face=}\StringTok{"bold"}\NormalTok{)) }\SpecialCharTok{+}
    \FunctionTok{theme}\NormalTok{(}\AttributeTok{strip.text.x =} \FunctionTok{element\_text}\NormalTok{(}\AttributeTok{size =} \DecValTok{12}\NormalTok{)) }\SpecialCharTok{+} \FunctionTok{theme}\NormalTok{(}\AttributeTok{legend.position=}\StringTok{"none"}\NormalTok{)\}}

\FunctionTok{plot\_all\_obs}\NormalTok{(}\StringTok{"Age"}\NormalTok{, }\StringTok{"Weight"}\NormalTok{, dataset,  }
             \AttributeTok{transfx =} \StringTok{" (in days)"}\NormalTok{, }\AttributeTok{transfy =} \StringTok{" (in grams)"}\NormalTok{)}
\end{Highlighting}
\end{Shaded}

\begin{figure}

{\centering \includegraphics{Report_files/figure-latex/vis_all_obs-1} 

}

\caption{Children weight observations over time}\label{fig:vis_all_obs}
\end{figure}

To observe data imbalance we can create multiple plots, grouping data by
number of observations for single child (see Figure \ldots). This can be
done by adding a line
\texttt{facet\_wrap(\textasciitilde{}.data{[}{[}"NObs"{]}{]})} to the
definition of \texttt{plot\_all\_obs()} function above (here we created
a separate function with alterations). We can also calculate absolute
frequency for number of visits depending on gender as in Table 2.

\begin{Shaded}
\begin{Highlighting}[]
\FunctionTok{plot\_facet\_var}\NormalTok{(}\StringTok{"Age"}\NormalTok{, }\StringTok{"Weight"}\NormalTok{, }\AttributeTok{df=}\NormalTok{dataset, }\StringTok{"NObs"}\NormalTok{, }
               \AttributeTok{transfx =} \StringTok{" (in days)"}\NormalTok{, }\AttributeTok{transfy =} \StringTok{" (in grams)"}\NormalTok{)}
\end{Highlighting}
\end{Shaded}

\begin{figure}

{\centering \includegraphics{Report_files/figure-latex/vis_facet_obs-1} 

}

\caption{Original observations by number of visits in the clinic}\label{fig:vis_facet_obs}
\end{figure}

\begin{Shaded}
\begin{Highlighting}[]
\FunctionTok{kable}\NormalTok{(}\FunctionTok{table}\NormalTok{(dataset}\SpecialCharTok{$}\NormalTok{Gender, dataset}\SpecialCharTok{$}\NormalTok{NObs) }\SpecialCharTok{/} \FunctionTok{rep}\NormalTok{(}\DecValTok{1}\SpecialCharTok{:}\DecValTok{5}\NormalTok{, }\AttributeTok{each =} \DecValTok{2}\NormalTok{))}
\end{Highlighting}
\end{Shaded}

\begin{longtable}[]{@{}lrrrrr@{}}
\caption{Number of participants per number of
observations}\tabularnewline
\toprule\noalign{}
& 1 & 2 & 3 & 4 & 5 \\
\midrule\noalign{}
\endfirsthead
\toprule\noalign{}
& 1 & 2 & 3 & 4 & 5 \\
\midrule\noalign{}
\endhead
\bottomrule\noalign{}
\endlastfoot
Boy & 45 & 62 & 86 & 75 & 15 \\
Girl & 46 & 79 & 83 & 68 & 9 \\
\end{longtable}

Here several observations can be made. First of all, both number or
observations and measurement time points vary strongly over
participants, which justifies use of multilevel models to analyse this
dataset. Secondly, 91 participants were only observed once after birth,
which makes these observations not longitudinal at all. They will be
excluded for later. Thirdly, dataset is balanced w.r.t. gender
distribution. Lastly, the original scale of the variables produces quite
large numerical values, which may negatively affect the numerical
estimation of models later. Therefore we need to rescale both target
variable and covariates. For this we will first convert grams into
kilograms and days into years and then center rescaled variables by
subtracting sample mean:

\begin{Shaded}
\begin{Highlighting}[]
\CommentTok{\# Exclude the data with 1 observation only (not longitudinal data)}
\NormalTok{dataset }\OtherTok{\textless{}{-}}\NormalTok{ dataset[dataset}\SpecialCharTok{$}\NormalTok{NObs }\SpecialCharTok{\textgreater{}} \DecValTok{1}\NormalTok{,]}

\NormalTok{dataset}\SpecialCharTok{$}\NormalTok{r\_Weight }\OtherTok{\textless{}{-}}\NormalTok{ dataset}\SpecialCharTok{$}\NormalTok{Weight }\SpecialCharTok{/} \DecValTok{1000}
\NormalTok{dataset}\SpecialCharTok{$}\NormalTok{r\_Birthweight }\OtherTok{\textless{}{-}}\NormalTok{ dataset}\SpecialCharTok{$}\NormalTok{Birthweight }\SpecialCharTok{/} \DecValTok{1000}
\NormalTok{dataset}\SpecialCharTok{$}\NormalTok{r\_Age\_years }\OtherTok{\textless{}{-}}\NormalTok{ dataset}\SpecialCharTok{$}\NormalTok{Age }\SpecialCharTok{/} \DecValTok{365}

\NormalTok{dataset}\SpecialCharTok{$}\NormalTok{c\_Age\_years }\OtherTok{\textless{}{-}}\NormalTok{ (dataset}\SpecialCharTok{$}\NormalTok{r\_Age\_years }\SpecialCharTok{{-}} \FunctionTok{mean}\NormalTok{(dataset}\SpecialCharTok{$}\NormalTok{r\_Age\_years))}
\NormalTok{dataset}\SpecialCharTok{$}\NormalTok{c\_Weight\_kg }\OtherTok{\textless{}{-}}\NormalTok{ (dataset}\SpecialCharTok{$}\NormalTok{r\_Weight }\SpecialCharTok{{-}} \FunctionTok{mean}\NormalTok{(dataset}\SpecialCharTok{$}\NormalTok{r\_Weight))}
\NormalTok{dataset}\SpecialCharTok{$}\NormalTok{c\_Birthweight\_kg }\OtherTok{\textless{}{-}}\NormalTok{ (dataset}\SpecialCharTok{$}\NormalTok{r\_Birthweight }\SpecialCharTok{{-}} \FunctionTok{mean}\NormalTok{(dataset}\SpecialCharTok{$}\NormalTok{r\_Birthweight))}

\NormalTok{df\_final }\OtherTok{=}\NormalTok{ dataset[}\FunctionTok{c}\NormalTok{(}\StringTok{"ChildID"}\NormalTok{, }\StringTok{"Gender"}\NormalTok{, }\StringTok{"NObs"}\NormalTok{, }\StringTok{"GenderID"}\NormalTok{,}
                      \StringTok{"c\_Weight\_kg"}\NormalTok{, }\StringTok{"c\_Birthweight\_kg"}\NormalTok{, }\StringTok{"c\_Age\_years"}\NormalTok{)]}

\FunctionTok{colnames}\NormalTok{(df\_final) }\OtherTok{\textless{}{-}} \FunctionTok{c}\NormalTok{(}\StringTok{"ChildID"}\NormalTok{, }\StringTok{"Gender"}\NormalTok{, }\StringTok{"NObs"}\NormalTok{, }\StringTok{"GenderID"}\NormalTok{,}
                       \StringTok{"Weight"}\NormalTok{, }\StringTok{"Birthweight"}\NormalTok{, }\StringTok{"Age"}\NormalTok{)}
\end{Highlighting}
\end{Shaded}

We can also observe, that intercept and slope varies from participant to
participant, which can also be modeled by multilevel models. In
addition, growth rate is not constant - babies grow very fast directly
after birth and then growth rate reduces after some time, - which means
that we can include some non-linearity in the model to improve the fit.

\section{Multilevel models of growing
flexibility}\label{multilevel-models-of-growing-flexibility}

We will fit different multilevel models of increasing complexity on this
and compare them based on different criteria. Using these models we can
try to answer following research questions, for example:

\begin{enumerate}
\def\labelenumi{\arabic{enumi}.}
\tightlist
\item
  How children's weight changes in the early stages of life?
\item
  Whether and how those growth patterns differ between individuals?
\item
  Are there any differences in development between children of different
  gender?
\item
  Is there any influence of birth weight on post-natal development?
\end{enumerate}

\subsection{Random intercept model}\label{random-intercept-model}

Multilevel models are sometimes called mixed models, and these two names
highlight different aspects of the very same method. They are called
``mixed'' because both fixed and random effects of predicted are
allowed, and ``multilevel'' stands for separate modeling of parameters
and their (co-)variances, i.e.~on different levels. Key idea is that the
observed variability comes from two sources: from both between-person
differences and within-person fluctuations. All multilevel models, as
stated above, assume observations from different subjects to be
independent, but correlated within the subject. Here the dependent
variable is assumed to be normally distributed.

The simplest multilevel model would be a random intercept model. That
means that we have constant predictions for a subject over time, but we
allow the intercept to vary over subjects. It can be written as follows
in multilevel notation:

\emph{Level 1:}
\(y_{ti} = \beta_{0i} + e_{ti}, \qquad e_{ti} \sim \mathcal{N}(0, \sigma_e^2)\)

\emph{Level 2:}
\(\beta_{0i} = \gamma_{00} + U_{0i}, \qquad U_{0i} \sim \mathcal{N}(0, \tau_{U_0}^2)\),
where:

\begin{itemize}
\tightlist
\item
  \(y_{ti}\) --- observation of dependent variable of unit \(i\) at time
  point \(t \in \{t_1, \ldots, t_N\}\);
\item
  \(\gamma_{00}\) --- fixed intercept;
\item
  \(U_{0i}\) --- random intercept, varies over individuals.
\end{itemize}

This model has 3 parameters to estimate: the fixed intercept
\(\gamma_{00}\), variance of random intercept \(\tau_{U_0}^2\) and
residual variance \(\sigma_e^2\).

We can fit a model using \texttt{lmer()} function from \texttt{lme4}
package (loaded above). The random effects are specified in brackets
with the grouping variable for the corresponding level behind the
vertical bar ``\textbar{}'':

\begin{Shaded}
\begin{Highlighting}[]
\NormalTok{model01 }\OtherTok{\textless{}{-}} \FunctionTok{lmer}\NormalTok{(Weight }\SpecialCharTok{\textasciitilde{}} \DecValTok{1} \SpecialCharTok{+}\NormalTok{ (}\DecValTok{1}\SpecialCharTok{|}\NormalTok{ChildID), }
                \AttributeTok{data=}\NormalTok{df\_final, }\AttributeTok{REML=}\ConstantTok{TRUE}\NormalTok{)}
\end{Highlighting}
\end{Shaded}

Model parameters are estimated using Maximum-Likelihood (ML) approach,
and in \texttt{lme4} package both unrestricted and restricted version
are implemented. Restricted ML (\texttt{REML\ =\ TRUE}) accounts for
uncertainty from simultaneous estimation of fixed parameters and
therefore delivers unbiased parameter variance estimates
\autocite{hoffman_2015}.

\begin{Shaded}
\begin{Highlighting}[]
\NormalTok{time\_grid }\OtherTok{\textless{}{-}} \FunctionTok{seq}\NormalTok{(}\FunctionTok{min}\NormalTok{(df\_final[}\StringTok{"Age"}\NormalTok{]), }\FunctionTok{max}\NormalTok{(df\_final[}\StringTok{"Age"}\NormalTok{]), }\AttributeTok{length.out =} \DecValTok{50}\NormalTok{)}

\NormalTok{plot\_results }\OtherTok{\textless{}{-}} \ControlFlowTok{function}\NormalTok{(data, model, model\_name,}
                         \AttributeTok{n =} \DecValTok{10}\NormalTok{, }\AttributeTok{ncols =} \DecValTok{5}\NormalTok{)\{}
  \FunctionTok{set.seed}\NormalTok{(}\DecValTok{42}\NormalTok{)}
\NormalTok{  sample\_n\_id }\OtherTok{\textless{}{-}} \FunctionTok{sample}\NormalTok{(}\FunctionTok{as.vector}\NormalTok{(}\FunctionTok{unique}\NormalTok{(data}\SpecialCharTok{$}\NormalTok{ChildID[}\FunctionTok{order}\NormalTok{(data}\SpecialCharTok{$}\NormalTok{ChildID)])), n)}
\NormalTok{  sample\_n\_id\_idx }\OtherTok{\textless{}{-}} \FunctionTok{which}\NormalTok{(data}\SpecialCharTok{$}\NormalTok{ChildID }\SpecialCharTok{\%in\%}\NormalTok{ sample\_n\_id)}
\NormalTok{  data\_n }\OtherTok{\textless{}{-}}\NormalTok{ data[sample\_n\_id\_idx, }
                 \FunctionTok{c}\NormalTok{(}\StringTok{"ChildID"}\NormalTok{, }\StringTok{"Birthweight"}\NormalTok{, }\StringTok{"GenderID"}\NormalTok{, }\StringTok{"Age"}\NormalTok{, }\StringTok{"Weight"}\NormalTok{)]}
\NormalTok{  unique\_data\_n }\OtherTok{\textless{}{-}} \FunctionTok{unique}\NormalTok{(data\_n[}\FunctionTok{c}\NormalTok{(}\StringTok{"ChildID"}\NormalTok{, }\StringTok{"Birthweight"}\NormalTok{, }\StringTok{"GenderID"}\NormalTok{)])}
\NormalTok{  newdata }\OtherTok{\textless{}{-}} \FunctionTok{cbind}\NormalTok{(unique\_data\_n, }
                   \FunctionTok{rep}\NormalTok{(}\FunctionTok{rownames}\NormalTok{(unique\_data\_n), }\AttributeTok{each =} \FunctionTok{length}\NormalTok{(time\_grid)))[,}\DecValTok{1}\SpecialCharTok{:}\DecValTok{3}\NormalTok{]}
\NormalTok{  newdata}\SpecialCharTok{$}\NormalTok{Age }\OtherTok{\textless{}{-}} \FunctionTok{rep}\NormalTok{(time\_grid, }\AttributeTok{each =}\NormalTok{ n)}
  
\NormalTok{  newdata}\SpecialCharTok{$}\NormalTok{Predictions }\OtherTok{\textless{}{-}} \FunctionTok{predict}\NormalTok{(model, newdata)}
  
\NormalTok{  g }\OtherTok{=} \FunctionTok{ggplot}\NormalTok{(data\_n, }\FunctionTok{aes}\NormalTok{(}\AttributeTok{x =}\NormalTok{ Age, }\AttributeTok{y =}\NormalTok{ Weight)) }\SpecialCharTok{+}
    \FunctionTok{geom\_point}\NormalTok{(}\AttributeTok{fill=}\StringTok{"grey"}\NormalTok{, }\AttributeTok{pch=}\DecValTok{21}\NormalTok{, }\AttributeTok{size=}\DecValTok{2}\NormalTok{, }\AttributeTok{stroke=}\FloatTok{1.25}\NormalTok{) }\SpecialCharTok{+}
    \FunctionTok{geom\_line}\NormalTok{(}\FunctionTok{aes}\NormalTok{(}\AttributeTok{group=}\NormalTok{ChildID)) }\SpecialCharTok{+} 
    \FunctionTok{facet\_wrap}\NormalTok{(}\SpecialCharTok{\textasciitilde{}}\NormalTok{ChildID, }\AttributeTok{ncol=}\NormalTok{ncols)}\SpecialCharTok{+}
    \FunctionTok{scale\_x\_continuous}\NormalTok{(}\AttributeTok{name =} \StringTok{"Age"}\NormalTok{) }\SpecialCharTok{+} 
    \FunctionTok{scale\_y\_continuous}\NormalTok{(}\AttributeTok{name =} \StringTok{"Weight"}\NormalTok{) }\SpecialCharTok{+}
    \FunctionTok{geom\_line}\NormalTok{(}\AttributeTok{data =}\NormalTok{ newdata, }
              \FunctionTok{aes}\NormalTok{(}\AttributeTok{x =}\NormalTok{ Age, }\AttributeTok{y =}\NormalTok{ Predictions, }\AttributeTok{group =}\NormalTok{ ChildID, }
                  \AttributeTok{color =} \FunctionTok{factor}\NormalTok{(ChildID)), }\AttributeTok{lwd=}\FloatTok{1.5}\NormalTok{) }\SpecialCharTok{+}
    \FunctionTok{ggtitle}\NormalTok{(model\_name)}
  
  \FunctionTok{return}\NormalTok{(g)\}}

\FunctionTok{plot\_results}\NormalTok{(df\_final, model01, }\AttributeTok{n=}\DecValTok{5}\NormalTok{, }\AttributeTok{ncols=}\DecValTok{5}\NormalTok{,}
             \AttributeTok{model\_name =} \StringTok{"Model 1: random intercept model"}\NormalTok{)}
\end{Highlighting}
\end{Shaded}

\begin{figure}

{\centering \includegraphics{Report_files/figure-latex/vis_model_fit-1} 

}

\caption{...}\label{fig:vis_model_fit}
\end{figure}

First we observe the model fit on a small subsample of participants. For
that we generate multiple time points and use fitted parameter values to
plot the model prediction.

In the Figure \ldots{} we can see that the model clearly underfits the
observations, which is quite natural: body weight increases while the
child grows. Another point is that fit of the model is singular,
i.e.~variance of random effects is estimated close to 0, which is not
very different from just fixed intercept. However, this model serves us
later as baseline.

\subsection{Random intercept, fixed linear time
model}\label{random-intercept-fixed-linear-time-model}

We can improve the model fit by adding covariates to the model, and a
first candidate would be a fixed linear effect of time \(\gamma_{10}\).
A model can be written as:

\emph{Level 1:}
\(y_{ti} = \beta_{0i} + \beta_{1i} \cdot x_{ti}+ e_{ti}, \quad e_{ti} \sim \mathcal{N}(0, \sigma_e^2)\)

\emph{Level 2:}
\(\beta_{0i} = \gamma_{00} + U_{0i}, \, \mbox{where} \;\; U_{0i} \sim \mathcal{N}(0, \tau_{U_0}^2), \quad \beta_{1i} =  \gamma_{10} \,,\)

where \(x_{ti}\) stands for time variable at observation \(t\) of
individual \(i\).

Hence, only the model for means is changed here and model for variance
remains the same. The resulting model will produce growth lines with
different intercepts but same slope for all subjects. We can now fit the
model and visualize the results:

\begin{Shaded}
\begin{Highlighting}[]
\NormalTok{model02 }\OtherTok{\textless{}{-}} \FunctionTok{lmer}\NormalTok{(Weight }\SpecialCharTok{\textasciitilde{}} \DecValTok{1} \SpecialCharTok{+}\NormalTok{ Age }\SpecialCharTok{+}\NormalTok{ (}\DecValTok{1}\SpecialCharTok{|}\NormalTok{ChildID), }
                    \AttributeTok{data=}\NormalTok{df\_final, }\AttributeTok{REML=}\ConstantTok{TRUE}\NormalTok{)}

\FunctionTok{plot\_results}\NormalTok{(df\_final, model02, }\AttributeTok{n=}\DecValTok{5}\NormalTok{, }\AttributeTok{ncol =} \DecValTok{5}\NormalTok{,}
             \AttributeTok{model\_name =} \StringTok{"Model 2: random intercept, fixed linear time"}\NormalTok{)}
\end{Highlighting}
\end{Shaded}

\begin{figure}

{\centering \includegraphics{Report_files/figure-latex/rint_fslp_model-1} 

}

\caption{Model predictions with random intercept and fixed slope}\label{fig:rint_fslp_model}
\end{figure}

Just after visual examination (see Figure 2) we can say that the model
fit improved greatly, even though the lines don't overlap with
observations perfectly. We can also review model summary for estimation
details. First of all, residuals are distributed around zero which may
be an indicator for a good fit. Random intercept \(U_{0i}\) has
estimated variance of \(\hat \tau_{U_0}^2 = 0.53\pm 0.73\) and residual
variance is estimated at \(\hat \sigma_e^2 = 1.76 \pm 1.33\).

Linear effect of time is estimated at \(\hat \gamma_{10} = 3.67\) is
considered significant (according to Wald test with
\(\frac{\hat\gamma_j}{SE(\hat\gamma_j)} \sim t_{N-f}\), where
\(j = 1,\ldots,f\) and \(f\) is number of fixed effects), however, the
fixed intercept (\(\hat \gamma_{00} = -0.02\)) is not significant. Small
correlation of fixed effects is also a good sign, meaning that fixed
parameters are uncorrelated.

\begin{Shaded}
\begin{Highlighting}[]
\FunctionTok{summary}\NormalTok{(model02)}

\CommentTok{\#Linear mixed model fit by REML. t{-}tests use Satterthwaite\textquotesingle{}s method }
\CommentTok{\#[lmerModLmerTest]}
\CommentTok{\# Formula: Weight \textasciitilde{} 1 + Age + (1 | ChildID)}
\CommentTok{\#    Data: df\_final}
\CommentTok{\# }
\CommentTok{\# REML criterion at convergence: 5358.1}
\CommentTok{\# }
\CommentTok{\# Scaled residuals: }
\CommentTok{\#     Min      1Q  Median      3Q     Max }
\CommentTok{\# {-}2.3687 {-}0.7542  0.0667  0.6669  4.4676 }
\CommentTok{\# }
\CommentTok{\# Random effects:}
\CommentTok{\#  Groups   Name        Variance Std.Dev.}
\CommentTok{\#  ChildID  (Intercept) 0.5329   0.730   }
\CommentTok{\#  Residual             1.7598   1.327   }
\CommentTok{\# Number of obs: 1481, groups:  ChildID, 477}
\CommentTok{\# }
\CommentTok{\# Fixed effects:}
\CommentTok{\#               Estimate Std. Error         df t value Pr(\textgreater{}|t|)    }
\CommentTok{\# (Intercept)   {-}0.01564    0.04851  456.71250  {-}0.322    0.747    }
\CommentTok{\# Age            3.67481    0.04495 1139.18833  81.746   \textless{}2e{-}16 ***}
\CommentTok{\# {-}{-}{-}}
\CommentTok{\# Signif. codes:  0 ‘***’ 0.001 ‘**’ 0.01 ‘*’ 0.05 ‘.’ 0.1 ‘ ’ 1}
\CommentTok{\# {-}{-}{-}{-}{-}{-}{-}{-}{-}{-}{-}{-}{-}{-}{-}{-}{-}{-}{-}{-}{-}{-}{-}{-}{-}{-}{-}{-}}
\CommentTok{\# Correlation of Fixed Effects:}
\CommentTok{\#     (Intr)}
\CommentTok{\# Age 0.002 }
\end{Highlighting}
\end{Shaded}

\subsection{Random linear time model}\label{random-linear-time-model}

We can go further and add random effect of time into the model. This
change will now affect model for the variance:

\emph{Level 1:}
\(y_{ti} = \beta_{0i} + \beta_{1i} \cdot x_{ti}+ e_{ti}, \quad e_{ti} \sim \mathcal{N}(0, \sigma_e^2)\)

\emph{Level 2:}
\(\beta_{0i} = \gamma_{00} + U_{0i}, \quad \beta_{1i} =  \gamma_{10} + U_{1i},\)

where \(U_{0i} \sim \mathcal{N}(0, \tau_{U_0}^2)\) and
\(U_{1i} \sim \mathcal{N}(0, \tau_{U_1}^2)\). This model has 6
parameters to estimate: fixed intercept \(\gamma_{00}\), fixed slope
\(\gamma_{10}\), residual variance \(\sigma_e^2\), variances of random
effects \(\tau_{U_0}^2\) and \(\tau_{U_1}^2\) as well as their
covariance \(\tau_{U_{01}}\).

According to \autocite{hoffman_2015} it is often useful to summarize
them as covariance matrix \(G\): \[G =  \begin{pmatrix}
                    \tau_{U_0}^2 & \tau_{U_{01}} \\ 
                    \tau_{U_{01}} & \tau_{U_1}^2
        \end{pmatrix},\] while residual variance is denoted as matrix R.
All variances and covariances are assumed constant over subjects.

We can then reformulate the model equation in vector notation, which
will result in \[
Y_i = X_i \cdot \mathbf{\gamma} + Z_i \cdot U_i + E_i,
\] where \(Y_i \in \mathbb{R}^{n_i}\) is target vector of individual
\(i\) (at \(n_i\) different time points), \(U_i = (U_{0i}, U_{0i})^T\)
is matrix of random effects and \(E_i\in \mathbb{R}^{n_i}\) summarizes
the residuals.

Here \(E_{i} \sim \mathcal{N}(0, R), \, R = \sigma_e^2 I_{n_i}\) is
assumed, however, another covariance structures (such as AR(1)) can be
explored and assessed.

Two design matrices \(X_i\) and \(Z_i\) are constructed as follows:

\[
X_i = \begin{pmatrix}
                 1 & t_1 \\
                 \vdots & \vdots \\
                 1 & t_{n_i}
              \end{pmatrix} \in \mathbb{R}^{n_i \times f} \quad
Z_i = \begin{pmatrix}
                 1 & t_1 \\
                 \vdots & \vdots \\
                 1 & t_{n_i}
              \end{pmatrix} \in \mathbb{R}^{n_i \times u},
\] where \(f\) is number of fixed effects, and \(u\) - number of random
effects in the model. Now, using vector notation we can also derive the
expression for the total variance (for an individual \(i\)): \[
V_i = \mathbb{V}ar(Y_{i}) = \mathbb{V}ar(X_i \cdot \mathbf{\gamma} + Z_i \cdot U_i + E_i) = \,Z_i^T \, \mathbb{V}ar(U_i)\,Z_i + R =  Z_i^T G Z_i + R.
\] Thanks to the construction of the design matrix \(Z_i\), the model
can accommodate data from irregular observations with different number
of time \(n_i\) points per person \(i\). Complete \(Z\) matrix is a
block-diagonal matrix with matrices \(Z_i, i = 1, \ldots, N\) on its
diagonal.

Another assumptions of the model is that no covariances are allowed for
the variance terms across levels, i.e.~between \(\tau_{U_0}^2\) and
\(\sigma_e^2\), and \(\tau_{U_1}^2\) and \(\sigma_e^2\) respectively.
Also, all units are assumed conditionally independent, meaning no
relationships within each of the \(U_{0i}, U_{1i}\), or \(e_{ti}\)
values across persons are allowed after including the model predictors.

Since we assume residuals to be multivariate normally distributed,
likelihood function for the model can be written as follows: \[
L = \prod_{i=1}^N (2\pi)^{-\frac{n}{2}} \cdot |V_i|^{-\frac{1}{2}} \cdot \exp{\left[{-\frac{1}{2}} (Y_i - X_i \mathbf{\gamma})^T V_i^{-1}(Y_i - X_i \mathbf{\gamma})\right]}  
\] with log-likelihood: \[
LL = -\frac{T}{2} \log(2\pi) -\frac{1}{2}\sum_{i=1}^N\log|V_i| - \frac{1}{2}\sum_{i=1}^N (Y_i - X_i \mathbf{\gamma})^T V_i^{-1}(Y_i - X_i \mathbf{\gamma})
\] Newton-Raphson algorithm is then used update estimates of matrix
\(V_i\) iteratively using derivatives of LL-function. For restricted ML
method following function is optimized: \[
LL_{\mbox{REML}} = -\frac{T-f}{2} \log(2\pi) -\frac{1}{2}\sum_{i=1}^N\log|V_i| - \frac{1}{2}\sum_{i=1}^N (Y_i - X_i \mathbf{\gamma})^T V_i^{-1}(Y_i - X_i \mathbf{\gamma}) - \frac{1}{2}\sum_{i=1}^N\log|X_i^T  V_i^{-1} X_i|,
\] for model with \(f\) fixed effects and \(N\) observations in total.

Random effects \(U_i\) are not included in the likelihood, therefore
they can be predicted with \(V_i\), \(X_i\) and \(Y_i\) given:
\(U_i = G_i Z_i^T V_i^{-1} (Y_i - X_i\mathbf\gamma)\). Fixed effects
\(\mathbf\gamma\) and their SEs also can be estimated with \(V_i\)'s
given (similar to GLM estimation): \[
\mathbf\gamma = \left(\sum_{i=1}^N(X_i^T V_i^{-1} X_i)\right)^{-1}\sum_{i=1}^N(X_i^T V_i^{-1} Y_i) \quad \mbox{and} \quad 
\mathbb{C}ov(\mathbf{\gamma}) = \left(\sum_{i=1}^N(X_i^T V_i^{-1} X_i)\right)^{-1}
\]\\
To fit a random linear time model we include \texttt{Age} variable as
random factor in addition to the intercept:

\begin{Shaded}
\begin{Highlighting}[]
\NormalTok{model03 }\OtherTok{\textless{}{-}} \FunctionTok{lmer}\NormalTok{(Weight }\SpecialCharTok{\textasciitilde{}} \DecValTok{1} \SpecialCharTok{+}\NormalTok{ Age }\SpecialCharTok{+}\NormalTok{ (}\DecValTok{1} \SpecialCharTok{+}\NormalTok{ Age}\SpecialCharTok{|}\NormalTok{ChildID), }
                    \AttributeTok{data=}\NormalTok{df\_final, }\AttributeTok{REML=}\ConstantTok{TRUE}\NormalTok{)}
\FunctionTok{plot\_results}\NormalTok{(df\_final, model03, }\AttributeTok{n=}\DecValTok{5}\NormalTok{, }\AttributeTok{ncols=}\DecValTok{5}\NormalTok{,}
             \AttributeTok{model\_name =} \StringTok{"Model 3: random linear time (RLT)"}\NormalTok{)}
\end{Highlighting}
\end{Shaded}

\begin{figure}

{\centering \includegraphics{Report_files/figure-latex/rint_rslp_model-1} 

}

\caption{Model predictions with random intercept and random slope}\label{fig:rint_rslp_model}
\end{figure}

Here we can see slight differences between the slopes for different
subjects. However, the model fit is singular again, here both random
effects are perfectly correlated. Still, we can compare three existing
models using \texttt{anova()} function from \texttt{lmerTest} package
for likelihood-ratio test. Results are collected in the Table 3, and we
can see that inclusion of time variable and even its random effect
increased the quality of the model.

\begin{Shaded}
\begin{Highlighting}[]
\FunctionTok{anova}\NormalTok{(model01, model02, model03)}
\end{Highlighting}
\end{Shaded}

\begin{longtable}[]{@{}lrrrrrrl@{}}
\caption{Unconditional model comparison}\tabularnewline
\toprule\noalign{}
Model & \# par & AIC & BIC & -2*LL & \(\chi^2_d\) & d & p-value \\
\midrule\noalign{}
\endfirsthead
\toprule\noalign{}
Model & \# par & AIC & BIC & -2*LL & \(\chi^2_d\) & d & p-value \\
\midrule\noalign{}
\endhead
\bottomrule\noalign{}
\endlastfoot
model01 & 3 & 7708.2 & 7724.1 & 7702.2 & & & \\
model02 & 4 & 5357.5 & 5378.7 & 5349.5 & 2352.620 & 1 & \textless{}
2.2e-16 \\
model03 & 6 & 5286.2 & 5318.0 & 5274.2 & 75.362 & 2 & \textless{}
2.2e-16 \\
\end{longtable}

\subsection{Conditional linear model}\label{conditional-linear-model}

We now explored different assumptions about the time structure. The
potential for further improvement remains, for example, through the
inclusion of additional covariates. Here we have two time-invariant
regressors - weight at birth as baseline observations (numeric) and
gender(binary), - that can be included in the model. Theoretically,
birth weight not only influences intercept, but also can have an impact
on growth rate after birth. For example, children that are born
premature, often have health issues directly after birth and may grow
slower. It is also interesting to study, whether differences between
development speed of boys and girls can be seen during early stages of
life.

Since it is possible to test all hypotheses at once, we will include
both fixed effects and also consider interactions between them and time
variable. Then the equation on the level 1 will be written as follows:
\[
y_{ti} = (\gamma_{00} + U_{0i}) + (\gamma_{10} + U_{1i}) \cdot x_{ti} + \gamma_{20} \cdot b_{ti} + \gamma_{30} \cdot g_{ti} + \gamma_{40} \cdot x_{ti}\cdot b_{ti} + \gamma_{50} \cdot x_{ti}\cdot g_{ti} + e_{ti}
\] where \(b_{ti}\) and \(g_{ti}\) denote birth weight and gender
respectively. Design matrix \(X_i\) is updated accordingly by two
columns and \(Z_i\) remains unchanged. To fit such model, we include
respective variables into the formula and then review model summary (see
Table 4).

\begin{Shaded}
\begin{Highlighting}[]
\NormalTok{model04 }\OtherTok{\textless{}{-}} \FunctionTok{lmer}\NormalTok{(Weight }\SpecialCharTok{\textasciitilde{}} \DecValTok{1} \SpecialCharTok{+}\NormalTok{ Age }\SpecialCharTok{*}\NormalTok{ Birthweight }\SpecialCharTok{+}\NormalTok{ Age }\SpecialCharTok{*}\NormalTok{ GenderID }\SpecialCharTok{+}\NormalTok{ (}\DecValTok{1}\SpecialCharTok{|}\NormalTok{ChildID), }
                \AttributeTok{data=}\NormalTok{df\_final, }\AttributeTok{REML=}\ConstantTok{TRUE}\NormalTok{)}
\end{Highlighting}
\end{Shaded}

\begin{Shaded}
\begin{Highlighting}[]
\FunctionTok{summary}\NormalTok{(model04)}
\FunctionTok{anova}\NormalTok{(model02, model04)}
\end{Highlighting}
\end{Shaded}

Residuals are again distributed around zero. Random intercept \(U_{0i}\)
has estimated variance of \(\hat \tau_{U_0}^2 = 0.25 \pm 0.50\), which
has decreased compared to corresponding unconditional model (model02).
Residual variance is estimated at \(\hat \sigma_e^2 = 1.76 \pm 1.33\)
(same as in the second model). Fixed intercept estimate increased from
\(-0.015\) to \(0.23\), and fixed slope increased a little as well. Both
new fixed effects seem to be significantly larger/smaller than 0 as well
as interaction term between birth weight and age, however, absolute
value of estimated parameters are rather small, except for fixed slope
parameter. Also the test shows that the addition of these covariates
significantly improves the fit (see Table 5).

\begin{longtable}[]{@{}
  >{\raggedright\arraybackslash}p{(\columnwidth - 10\tabcolsep) * \real{0.3590}}
  >{\raggedleft\arraybackslash}p{(\columnwidth - 10\tabcolsep) * \real{0.1282}}
  >{\raggedleft\arraybackslash}p{(\columnwidth - 10\tabcolsep) * \real{0.1538}}
  >{\raggedleft\arraybackslash}p{(\columnwidth - 10\tabcolsep) * \real{0.1154}}
  >{\raggedleft\arraybackslash}p{(\columnwidth - 10\tabcolsep) * \real{0.1154}}
  >{\raggedright\arraybackslash}p{(\columnwidth - 10\tabcolsep) * \real{0.1282}}@{}}
\caption{Conditional model parameter estimation}\tabularnewline
\toprule\noalign{}
\begin{minipage}[b]{\linewidth}\raggedright
Predictor
\end{minipage} & \begin{minipage}[b]{\linewidth}\raggedleft
Estimate
\end{minipage} & \begin{minipage}[b]{\linewidth}\raggedleft
Std. Error
\end{minipage} & \begin{minipage}[b]{\linewidth}\raggedleft
df
\end{minipage} & \begin{minipage}[b]{\linewidth}\raggedleft
t value
\end{minipage} & \begin{minipage}[b]{\linewidth}\raggedright
p-value
\end{minipage} \\
\midrule\noalign{}
\endfirsthead
\toprule\noalign{}
\begin{minipage}[b]{\linewidth}\raggedright
Predictor
\end{minipage} & \begin{minipage}[b]{\linewidth}\raggedleft
Estimate
\end{minipage} & \begin{minipage}[b]{\linewidth}\raggedleft
Std. Error
\end{minipage} & \begin{minipage}[b]{\linewidth}\raggedleft
df
\end{minipage} & \begin{minipage}[b]{\linewidth}\raggedleft
t value
\end{minipage} & \begin{minipage}[b]{\linewidth}\raggedright
p-value
\end{minipage} \\
\midrule\noalign{}
\endhead
\bottomrule\noalign{}
\endlastfoot
(Intercept) & 0.23682 & 0.05842 & 429.06364 & 4.054 & 5.98e-05 \\
Age & 3.71449 & 0.06264 & 1168.76560 & 59.295 & \textless{} 2e-16 \\
Birthweight & 0.91200 & 0.08321 & 429.88095 & 10.960 & \textless{}
2e-16 \\
GenderID & -0.49817 & 0.08329 & 442.78809 & -5.981 & 4.58e-09 \\
Age \(\times\) Birthweight & 0.18107 & 0.08929 & 1183.33744 & 2.028 &
0.0428 \\
Age \(\times\) GenderID & -0.09151 & 0.08915 & 1172.60511 & -1.026 &
0.3049 \\
\end{longtable}

\begin{longtable}[]{@{}lrrrrrrl@{}}
\caption{Conditional vs.~unconditional model comparison}\tabularnewline
\toprule\noalign{}
Model & \# par & AIC & BIC & -2*LL & \(\chi^2_d\) & d & p-value \\
\midrule\noalign{}
\endfirsthead
\toprule\noalign{}
Model & \# par & AIC & BIC & -2*LL & \(\chi^2_d\) & d & p-value \\
\midrule\noalign{}
\endhead
\bottomrule\noalign{}
\endlastfoot
model02 & 4 & 5357.5 & 5378.7 & 5349.5 & & & \\
model04 & 8 & 5219.0 & 5261.4 & 5203.0 & 146.59 & 4 & \textless{}
2.2e-16 \\
\end{longtable}

\subsection{Beyond linearity: quadratic time
effect}\label{beyond-linearity-quadratic-time-effect}

However, there is still some possibility for improvement. During data
exploration we mentioned that the growth doesn't seem to be linear:
weight increases very quickly short after birth but then the growth
slows down. We can try to account for that in the model by including
non-linear transformation of the time variable into the models for mean
and/or for variance.

\begin{Shaded}
\begin{Highlighting}[]
\NormalTok{model05 }\OtherTok{\textless{}{-}} \FunctionTok{lmer}\NormalTok{(Weight }\SpecialCharTok{\textasciitilde{}} \DecValTok{1} \SpecialCharTok{+}\NormalTok{ Age }\SpecialCharTok{+} \FunctionTok{I}\NormalTok{(Age}\SpecialCharTok{\^{}}\DecValTok{2}\NormalTok{) }\SpecialCharTok{+}\NormalTok{ (}\DecValTok{1}\SpecialCharTok{|}\NormalTok{ChildID), }
                \AttributeTok{data=}\NormalTok{df\_final, }\AttributeTok{REML=}\ConstantTok{TRUE}\NormalTok{)}
\NormalTok{model06 }\OtherTok{\textless{}{-}} \FunctionTok{lmer}\NormalTok{(Weight }\SpecialCharTok{\textasciitilde{}} \DecValTok{1} \SpecialCharTok{+}\NormalTok{ Age }\SpecialCharTok{+} \FunctionTok{I}\NormalTok{(Age}\SpecialCharTok{\^{}}\DecValTok{2}\NormalTok{) }\SpecialCharTok{+}\NormalTok{ (}\DecValTok{1} \SpecialCharTok{+}\NormalTok{ Age}\SpecialCharTok{|}\NormalTok{ChildID), }
                    \AttributeTok{data=}\NormalTok{df\_final, }\AttributeTok{REML=}\ConstantTok{TRUE}\NormalTok{)}
\FunctionTok{plot\_results}\NormalTok{(df\_final, model06, }\AttributeTok{n=}\DecValTok{5}\NormalTok{, }\AttributeTok{ncols=}\DecValTok{5}\NormalTok{,}
             \AttributeTok{model\_name =} \StringTok{"Model 6: random intercept, random linear, fixed quadratic time"}\NormalTok{)}
\end{Highlighting}
\end{Shaded}

\begin{figure}

{\centering \includegraphics{Report_files/figure-latex/rint_fquad_model-1} 

}

\caption{Model predictions with quadratic time}\label{fig:rint_fquad_model}
\end{figure}

\begin{Shaded}
\begin{Highlighting}[]
\FunctionTok{anova}\NormalTok{(model02, model05, model06)}
\end{Highlighting}
\end{Shaded}

When extending the model to account for non-linear development, a
quadratic function of time yielded a more accurate representation of the
data (compare Table 6). In this case, models incorporating both fixed
and random linear time effects were appropriate, however, quadratic term
of time should be rather included as fixed effect only. This indicates
that the observed trajectories are not purely linear but exhibit
curvature over time. However, quadratic function of time cannot be an
optimal solution for long-term prediction. By definition it starts to
decrease after some time with the same rate it was increasing before,
which cannot be an adequate model for body weight.

\begin{longtable}[]{@{}lrrrrrrl@{}}
\caption{Model comparison for linear vs.~quadratic time}\tabularnewline
\toprule\noalign{}
Model & \# par & AIC & BIC & -2*LL & \(\chi^2_d\) & d & p-value \\
\midrule\noalign{}
\endfirsthead
\toprule\noalign{}
Model & \# par & AIC & BIC & -2*LL & \(\chi^2_d\) & d & p-value \\
\midrule\noalign{}
\endhead
\bottomrule\noalign{}
\endlastfoot
model02 & 4 & 5357.5 & 5378.7 & 5349.5 & & & \\
model05 & 5 & 4356.0 & 4382.5 & 4346.0 & 1003.54 & 1 & \textless{}
2.2e-16 \\
model06 & 7 & 4080.5 & 4117.6 & 4066.5 & 279.49 & 2 & \textless{}
2.2e-16 \\
\end{longtable}

\section{Summary and outlook}\label{summary-and-outlook}

Analyzing longitudinal data with multilevel models reveals some key
practical and methodological aspects. First key observation concerns the
robustness of multilevel models in the presence of missing and
imbalanced data. Longitudinal datasets are often characterized by
unequal numbers of observations across individuals or irregular
measurement intervals. The applied modeling framework proved capable of
handling such cases effectively, which represents a substantial
advantage over more traditional approaches such as repeated-measures
ANOVA. Furthermore, the results emphasize the crucial role of data
pre-processing. Proper handling of missing values and scaling of
variables are essential not only for ensuring model convergence but also
for enabling meaningful interpretation of the estimated parameters.

Several directions for future research can be formulated based on the
results of this analysis. Firstly, other structures for variance models
can be explored. Here, only unstructured variance model was used, but it
is possible that other structures such as compound symmetry (CS) or
autoregressive structure will describe the variability in the
observations better. Secondly, alternative non-linear transformations of
time variable can be explored. For example, logarithmic or exponential
functions may provide a better fit for certain types of growth or decay
processes in long term studies, particularly when changes occur more
rapidly at early time points and stabilize later on. The flexibility of
multilevel models allows such transformations to be incorporated
straightforwardly. Model comparison techniques and diagnostic tools
should be employed to assess whether such extensions are justified by
the data.

In conclusion, multilevel modeling provides a powerful and flexible
framework for analyzing longitudinal data, enabling researchers to
capture both overall trends and individual trajectories. Future
developments and extensions of this approach hold considerable potential
for advancing empirical research in this field.

\printbibliography

\end{document}
